% !TeX root = ../papers/nc-p3.tex
\section{各因子の接平面の次数と三次因子の特異点}
\label{sec:cubic-octic-tangent-planes}
\begin{theorem}\label{thm:cubic-octic-tangent-planes}
定理\ref{thm:low-degree-octic-spinor}の各既約因子 $S$ に対し、$e=\deg S$,
$m=\sum_{p\in X}\operatorname{mult}_p S$ とする。出力の一般三空間 $W\subset k^4$ で、
$X$ の外の $S$ 上で $N|_W$ が階数1となる点の長さは $m-6e$ である。
特に $m\ge6e$。可約な場合の三次因子が印で三重でなければ、それは正規であり、
$(n,b)=(15,3),(14,4),(15,4)$ に限られる。ここで $n$ は因子上の印の個数、$b$ はそのうちの二重点の個数である。
三つの二重の印は共線でない。
\end{theorem}
\begin{proof}[三行制限で数える残余十二点]
階数3の点で $N|_W$ の階数が1となるためには、$W$ は核を含み、核で割った正則三次元形式に対し
残る二平面が退化しなければならない。従って各入力点でこのような $W$ は滑らかな二次曲線をなす。
階数4の点ではそのような $W$ はない。一般 $W$ で $X$ 外の該当集合は有限であり、
因子同士の交叉や因子の特異曲線も避ける（それらの上の入射は二次元以下）。
$N|_W$ の六つの二次式は同時に零にならない。全零なら $N$ の階数は高々2で、
一般 $W$ は有限個の $y(p)^\perp$ と異なるからである。従って $\mathbb P^3\to\mathbb P^5$ の
二次写像によるVeronese平面の引戻しがこの有限スキームである。次数は $4\cdot2^3=32$ である。
例えば六つの式を三つの新しい一次座標の全積に等置すれば $\mathbb P^6$ の六二次完全交叉となり、
各点では単元の平方根の二枚で元のスキームを覆う。

各 $p\in X$ はこのスキームの単純点である。単元を消去した対称三行の一次部を $A(z)$ とすると、
$W$ を制限した二行の三成分は、残りの三空間内の一般二平面への $A(z)$ の制限である。
非零 $z$ の形式がそこで恒等零なら、階数2の $A(z)$ が与える二直線のいずれかにこの二平面が対応する。
そのような二平面は接三次上の高々一パラメータであるから一般には起こらない。従って三成分は独立な一次部を持つ。
印を除いた全長は12である。

点の吹上げでの $S$ の真変換を $S'$、核直線束の双対を $L$ とする。$W$ が核を含む条件は $L$ の
一般超平面曲線 $C$ を切る。$L^2=9e-m$, $HL=3e$ である。
$C$ 上の二束 $W/\mathcal O(-L)$ の形式の行列式は $4H-2L$ の切断となる。
一般 $W$ でその零スキームは各例外曲線との交叉を除いて先の入射の当該因子上の部分である。
$S'$ の特異点や射影の例外的な曲線でも、一般の超平面はその成分を含まず、二条件を同時に満たす点は上の次元勘定で除ける。
例外因子 $E$ の上では $N$ の階数が1なので、この行列式は $E|_C$ で消える。
実際元の二次小行列式は $\mathcal I_X$ に属し、局所的に二束の基底を持ち上げて計算すればよい。
$E\cdot C=m$（$2L=2(3H-E)|_{S'}$）であり、当該因子上の印の外の長さは
\[
                (4H-2L)L-m=m-6e
\]
以下となる。ここで一般超平面は例外曲線を成分に含まない。各例外平面上の余因子射は全二次式による閉埋込みで、
例外曲線上の $W$ の二条件が同時に起こる場所も有限である。従って行列式は $C$ の非零因子で、
交点数による上界を適用できる。もし核像が一点なら $C$ は空で、同じ数値式を用いる。
全因子にわたる上界の和は $60-6\cdot8=12$ に等しいため、全ての上界が等号である。
\end{proof}
\begin{proof}[三次因子の一般の二重の印]
三次因子 $S$ に三重の印がなければ $L=P$ であり $h^0(L)\le19-n$。
$m\ge18$ かつ $L^2=27-m>0$（定理\ref{thm:reducible-octic-factors}）より $n\le16$。
従って二重の印が二つ以上あり、二つの例外二次曲線の包から $h^0(L)\ge4$ となる。
故に $n\le15$、二重の印は三つ以上。孤立特異点の場合は高々四個、三点共線は不可能なので、
$m=n+b\ge18$ が表の三通りだけを残す。三次錐で頂点だけが特異ならこの範囲に入れない。

特異直線を持つ場合、全特異点はその直線に載るので $b=3,n=15$ である。この場合を排除する。
直線の横断座標を $u,v$、残りを $w,z$ とすると方程式は
$w f_2(u,v)+z g_2(u,v)+h_3(u,v)=0$。比 $[u:v]=[s:t]$ を加えると、各平面の残余は
$w f_2(s,t)+z g_2(s,t)+a h_3(s,t)=0$（$(u,v)=a(s,t)$）という直線となる。
係数に共通零点があれば元の式は可約なので、これは滑らかな射影直線束 $T=\mathbb F_j$ を与える。
$H=C_0+kf$ は基点なしで $2k-j=3$, $k\ge j\ge0$、従って $(j,k)=(1,2),(3,3)$。
$T$ から三次曲面への写像は特異直線の外で同型である。二重の印では接二次式が相異なる二直線に分かれ、
二項二次式 $w f_2+z g_2$ の二根に対応する二つの滑らかな局所枝がある。
それぞれの枝に三次の核座標を制限すると基底イデアルは単純である。
実際各枝の接方向の射影直線上で例外核写像は閉埋込みで次数1であり、一次部が枝の余接空間を張る。
頂点があっても印でないからそこに新たな基点は生じない。

前節の直交三束と射影束の構成を $T$ 上で用いると、$s=n+b=18$, $c_2(E)=0$、
四切断の共通零点は18個、一般三切断の長さも $2s-6H^2=18$ である。
$c_1(V)=a C_0+c f$ の自乗 $-ja^2+2ac$ は4の倍数で、$j$ は奇数ゆえ $a$ は偶数。
$Hf=1$ なので母線上は均等、捩りを除き $V$ は底から来る。$E$ は
$\mathcal O,\mathcal O(lf),\mathcal O(-lf)$ の和で、$0\le l\le H C_0=k-j$。
従って $H^1(E(-2H))=0$ である。実際 $K_T+2H=(2k-j-2)f=f$ で、
双対側は $\mathbb P^1$ 上の次数 $1,1\pm l$ の第一コホモロジーであり $l\le1$。
四番目の独立な切断は長さ18の完全交叉全体で消えることができず、矛盾する。
\end{proof}
